\chapter{CONCLUSIONS}\label{ch:conclusion}

This lab exercise successfully demonstrated the application of UML class diagramming for domain modeling in the context of the Multi-Tenant E-Commerce System. The key conclusions are:

\begin{itemize}
    \item \textbf{Domain Model Completeness:} A comprehensive domain model was developed encompassing \textbf{10 core classes and 37 total classes} that collectively represent all major business entities in the system.
    
    \item \textbf{Business Rule Integration:} Critical business constraints were identified and documented, ensuring that the conceptual model accurately reflects real-world operational policies.
    
    \item \textbf{Multi-Tenant Architecture Support:} The domain model explicitly incorporates multi-tenancy through vendor-specific ownership of domain entities, laying the groundwork for data isolation in implementation.
    
    \item \textbf{Tool Proficiency:} Lucidchart proved effective for creating professional UML diagrams with proper notation, multiplicities, and relationships.
\end{itemize}

\noindent The domain model created in this lab bridges the gap between requirements analysis (use cases) and system design, ensuring that all stakeholders share a common understanding of the business domain before implementation begins. This conceptual foundation is essential for building a robust, scalable, and maintainable multi-tenant e-commerce platform.